% Lösungen Einheit 3 von 5 – Punkte und Werte (zusammenbau v0.1)
\einheitenkopf{Einheit 3 von 5 · Punkte und Werte}

\erg{19}{a) $3 \cdot 4 - 2 = 10$}
\erg{20}{a) $f(4) = 2 \cdot 4 + 5 = 13$ \quad b) $f(-3) = 4 \cdot (-3) - 1 = -13$ \quad c) $f(1{,}5) = 4 \cdot 1{,}5 - 3 = 3$ \quad d) $11 = 2x + 3$, also $x = 4$ \quad e) Ja: $f(3) = 8$ \quad f) $0 = 3x - 12$, also $x = 4$ \quad g) Gerade durch $(0|6)$ und $(1|2)$; $0 = -4x + 6$, Nullstelle $x = 1{,}5$}
\erg{21}{a) $5$, $4$, $3$, $2$, $1$}
\erg{22}{a) $S(0|-8)$}
\erg{23}{a) Sara hat den y-Achsenabschnitt angegeben; richtig: $0 = 2x - 6$, also $x = 3$. \quad b) Nein: Eine steigende Gerade schneidet die x-Achse genau einmal. \quad c) $n = -1$ (Wert bei null), $m = 3$ (Zuwachs je Schritt): $f(x) = 3x - 1$ \quad d) $0 = -1{,}5x + 36$, $x = 24$: nach 24 Sekunden}
